\documentclass[10pt,a4paper]{amsart}
\usepackage[margin=19mm]{geometry}
\usepackage{amsmath,amssymb,mathtools}
\usepackage[scr=rsfs,cal=euler]{mathalfa}
\usepackage{xfrac}
\usepackage[unicode,hidelinks,bookmarks=false]{hyperref}
\newcommand{\ie}{i.e.\ }
\newcommand{\cf}{cf.\ }
\newcommand{\resp}{respectively\ }
\newcommand{\Subsection}{\subsection}
\providecommand{\nobreakdash}{\nobreak\hbox{-}}
\setlength{\emergencystretch}{2em}
\multlinegap0pt
\usepackage{amssymb}
\usepackage{bm}
\usepackage{float}
\usepackage{mathtools}
\newtheorem{lem}{Lemma}
\title{Gaps between quadratic forms}
\author{Siddharth Iyer}
\date{Source: arXiv:2505.23428v1 (CC BY 4.0)}
\begin{document}
\maketitle

\noindent\emph{Source and licence.} Gaps between quadratic forms. Version arXiv:2505.23428v1 (CC BY 4.0), distributed by arXiv under the Creative Commons Attribution 4.0 International licence, http://creativecommons.org/licenses/by/4.0/, as recorded in the arXiv licence metadata for this exact version and shown on the abstract page https://arxiv.org/abs/2505.23428v1. Full licence notice: \texttt{licenses/arxiv-2505.23428-CC-BY-4.0.txt}.\par\smallskip

\noindent\textit{Editorial scope.} Counted unit: the article's lem Twistedsumlem (source0.tex lines 578--584). The article's complete original proof of it is delivered with it (source0.tex lines 585--608, one \texttt{proof} environment of 24 lines). No mathematics is written, repaired or re-derived: the statement and the proof are the article's own lines, in the article's own notation, and no retained line is substituted or re-flowed.  Retained uncounted as the source-local context the proof needs: lines 550--557.  Nothing was omitted from any retained span, so the package carries no omission record.  No retained line contains a citation, so the wrapper carries no bibliography and no bibliography entry.  No retained line contains a figure environment, an image-inclusion command or drawing code, so no image is delivered and the package declares no asset dependency.  Editorial formatting only, in addition: an AMS \texttt{amsart} preamble that restates the article's own declarations and macros used by the retained lines, namely the environments \texttt{lem} on separate counters, together with the standard packages those lines need.  The retained environments print in this document as Lemma 1 in order of appearance, whereas the article prints the same items with the numbering of its own full document.  The complete native source of the article as distributed in the arXiv e-print for this exact version is bundled unchanged as \texttt{sources/178998/source0.tex}.
\begin{lem}
\label{lambdaconvbound}
We have
\begin{align*}
\overline{\lambda_{a}}(n) = \frac{f_{\lambda_{a}}(n)}{n},
\end{align*}
where $|f_{\lambda_{a}}(z)| \leq C_{a}$ whenever $z$ is an odd number and $C_{a}>0$ is some positive constant.
\end{lem}

\begin{lem}
\label{Twistedsumlem}
When $\rho_{b}$ is a non-trivial multiplicative character we have
\begin{align*}
J(\rho_{b},b,a,x) \ll x \log(x).
\end{align*}
\end{lem}

\begin{proof}
Let $1\leq \alpha \leq b$ be an integer with $(\alpha,b)=1$. We write 
\begin{align*}
I(\alpha,b,a,x)&:=\sum_{\substack{n \leq x \\ n \equiv \alpha \pmod b}}\hspace{-0.5cm}\eta_{a}(n) = \sum_{\substack{n \leq x \\ n \equiv \alpha \pmod b}}\hspace{-0.5cm}n \lambda_{a}(n)\\\\
&= \sum_{\substack{n \leq x \\ n \equiv \alpha \pmod b}}n \sum_{d|n}\overline{\lambda_{a}}(d).
\end{align*}
Using the notation and result of Lemma \ref{lambdaconvbound} we can continue the above equation as
\begin{align*}
&\sum_{\substack{n \leq x \\ n \equiv \alpha \pmod b}}\left( \sum_{md = n}m f_{\lambda_{a}}(d)\right)= \sum_{d \leq x, \ (d,b) = 1}f_{\lambda_{a}}(d)\sum_{\substack{md \leq x \\ md \equiv \alpha \pmod b}} m\\
&= \sum_{d \leq x, \ (d,b) = 1}f_{\lambda_{a}}(d)\left(\left(\sum_{\substack{md \leq x \\ md \equiv 0 \pmod b}} m\right) + O(x/d)\right)\\
&=\left(\sum_{d \leq x, \ (d,b) = 1}f_{\lambda_{a}}(d)\sum_{\substack{m \leq x/d \\ m \equiv 0 \pmod b}} m \right) + O(x\log(x)).
\end{align*}
Noting that
\begin{align*}
J(\rho_{b},b,a,x) = \sum_{\alpha = 1, \ (\alpha,b) = 1}^{b}\rho_{b}(\alpha)\cdot I(\alpha,b,a,x),
\end{align*}
we have
\begin{align*}
J(\rho_{b},b,a,x) &= \left(\sum_{\alpha = 1, \ (\alpha,b) = 1}^{b}\rho_{b}(\alpha)\right)\left(\sum_{d \leq x, \ (d,b) = 1}f_{\lambda_{a}}(d)\sum_{\substack{m \leq x/d \\ \ m \equiv 0 \pmod b}} m \right)\\
&\hspace{0.5cm}+ O(x\log(x))\\
&= O(x\log(x)).
\end{align*}
\par \vspace{-\baselineskip} \qedhere
\end{proof}

\end{document}
